\chapter{Conclusion}
\label{chap:conclusion}

This project asked whether external public data can be used to predict business needs and inform
commercial banking strategies. The answer found is a partial yes. Company filings can order UK
SMEs by their likelihood of taking a first Lloyds charge at nine times better than picking at
random. News coverage, at the depth that open sources are able to give, cannot.

\textbf{The data}. A Companies House pipeline was built on two routes, since the bulk product
carries the population cheaply but holds no charges, and the rate limit makes pulling the whole
register through the API impractical. Together they collected 839,295 companies, of which 567,835
passed the SME filter, with a history of 399,703 registered charges. The label was then matched
from the \textit{persons\_entitled} field, which is free text and holds 33,822 distinct lender
strings, using one regular expression for each Lloyds Banking Group brand. However, one mistake
found here is worth noting again. Matching ``bank of scotland'' alone also matches ``The Royal
Bank of Scotland'', which belongs to NatWest Group, and this quietly produced 14,305 wrong
charges, over a quarter of all matches, before it was corrected. With a cutoff of 1 January 2021
and a five-year window, the modelling table holds 360,600 prospects of which 1,203 converted, a
rate of 0.33 percent.

\textbf{The models}. Logistic regression, XGBoost and MLP were chosen for their difference in
explainability rather than for expected accuracy alone, and each was tuned with five-fold
cross-validation and grid search. The class imbalance was deliberately kept rather than
rebalanced, since forcing the dataset into balance would manipulate the model in an undesired way
and fail to reflect the real population.

\textbf{The metrics}. Precision@K and lift over the base rate are reported, and accuracy is not,
because a classifier answering ``no conversion'' for every company would already be correct
99.67 percent of the time. Average precision was used as the grid search metric, and this choice
was tested rather than assumed. Three grid searches were run on identical data and seeds, and the
model tuned directly on Precision@100 produced the worst Precision@K at every K, including
K = 100 itself, at 1.0 percent against 4.0 percent. Therefore, optimising for one list size does
not reliably produce the best result at that size.

\textbf{The filings}. Roughly a quarter of accounts are filed as scanned documents with no text
layer, which places them outside the machine-readable route entirely. A Qwen3-VL-8B model was
fine-tuned with QLoRA to read tables out of them, reaching a mean TEDS of 0.904 against 0.438 for
the Table Transformer and Tesseract approach, which struggled on the sparse and frameless tables
that financial accounts are made of.

\textbf{The media index}. Per-company news was tested first and abandoned, as SMEs almost never
appear in coverage and there is no reliable key to join an article to a company number. Therefore,
the signal was moved up to a weekly index of 21 SIC sections by twelve UK regions, each cell
z-scored against its own 52-week baseline. It was pulled with a single aggregated query rather
than a loop over cells, and every query is dry-run against a cost cap first, which kept the whole
pull inside the monthly free tier.

\textbf{What works and what does not}. The pipeline, the label and the ranking all work. The
shipped model reaches Precision@100 of 3 percent against a base rate of 0.33 percent, a lift of
nine times, and the same lift is held through K = 500. Since K is a business parameter and not a
statistical one, holding the lift as the list grows matters more than the very top of it alone.
Logistic regression is the model shipped, as the three sit inside the noise at this number of
positives, and explainability therefore decided it, because the coefficients are what let a lead
be justified to a relationship manager. The media index does not contribute. Tone gives a lift of
$1.00\times$, sector-only volume gives $p = 0.86$ once deseasonalised, and adding both as
features moves Precision@K by approximately zero at every K. Objective five was to test exactly
this, and it is therefore completed, only with a negative answer. The index is kept as a soft
re-ordering at a weight of 0.15 on real cells only, and it never removes a company from the list.
The structural features, however, remain current values taken from the register rather than
values reconstructed at the cutoff, and this is the main limitation carried into the result.

\textbf{Further work}. The filing history is already pulled and sits unused. Rebuilding the
account category as it stood at the cutoff would remove the largest remaining leakage in the
feature table, and it needs no new data collection. The features read out of the filings are
extracted but not yet joined into the modelling table, so the model still runs with no measure of
financial size, and tangible fixed assets are the closest thing the public record has to a
collateral base. Lastly, the most interesting outcome of this project was not expected. Volume at
sector by region level shows a strong association of $z = +6.77$ on real cells, and yet the
hottest quintile converts at $0.97\times$ the baseline rate, which is marginally below chance.
Something is being detected. However, it does not point the way the hypothesis assumed. One
explanation worth testing is that funding coverage of a sector rises when that sector is under
stress, which would mix distress and expansion together inside the same count, and separating
them would form a useful study on its own.
